% ============================================================
% Chapter 3: System Architecture and PCB Design
% ============================================================
\chapter{System Architecture and PCB Design}
\label{chap:architecture}

This chapter details the technical implementation of the retrofitted test-diagnosis module, translating the design principles discussed in~\Cref{sec:design_principles} into a concrete hardware architecture. The system employs a modular design across four distinct~\glspl{pcb} to enhance isolation, testability, and maintainability, reflecting common practice for complex electronic systems. Detailed schematics and layout diagrams are provided in~\Cref{app:schematics}.

% ------------------------------------------------------------
\section{System Architecture Overview}
\label{sec:overall_architecture}

\begin{figure}[htbp]
	\centering
	% High-level block diagram of the retrofitted test-diagnosis system.
% Kept as diagram source rather than a bitmap so that it diffs, reviews and
% rebuilds deterministically with the rest of the document. Derived from the
% block diagram of the published paper; the RTC node names the STM32's internal
% RTC with its 32.768 kHz crystal, which is what the Rev 1.0 schematic and layout
% show. The board assignment follows the Rev 1.0 Eagle layouts: the BOD
% supervisor (Y3) and the CD4050 level shifters for the X4 inputs sit on the Base
% Board, the watchdog (Y4, Y5, Q1) and the debouncing (Y6) on the Upper Board.
% The BOD output drives the watchdog's reset transistor Q1 over X9/X8, so its
% reset reaches the FPGA through the Upper Board. The outputs to the system
% connector X4 leave the Base Board through the 54LS245-type drivers Y8 and Y12
% (schematic sheets 2 and 3); their arrow leaves the board below the LED Board.
% Edge labels sit on the side of their arrow away from the neighbouring label,
% and the STM32 peripherals stand far enough from the STM32 for a visible arrow.
% The 5 V supply and ground of the Base Board reach the Diagnostics Board through
% X13/X12 (schematic sheet 4); on the Base Board side of the barrier the 5 V (net
% PWR_SUPPLY_5V) feeds only the suppressor diode D16 and the input of the isolated
% converter Y23 (the footprint Y22 is unfitted), drawn as the dark arrow. The LEDs of
% the optocouplers share only the ground (GND3) and are driven by the logger and
% reset lines through R82-R85 (X12 pins 2-5); the branch is their ground return.
% The FPGA reaches the NVRAM over address, data and enable lines (atop.vhd, Rev 1.0
% schematic: A0-A5, DQ0-DQ7, CE, OE, WE, BLE); the fitted NVRAM needs no STORE and
% RECALL, and the X4 STORE/RECALL inputs feed only the logger.
\resizebox{\textwidth}{!}{%
\begin{tikzpicture}[
    module/.style={draw, thick, rounded corners=3pt, minimum width=2.6cm,
                   minimum height=1.2cm, align=center, font=\small},
    fpga/.style={draw, very thick, rounded corners=5pt, minimum width=5.6cm,
                 minimum height=4cm, align=center, font=\normalsize\bfseries,
                 fill=gray!5},
    pcb/.style={draw, thick, rounded corners=8pt, inner sep=14pt},
    pcblabel/.style={font=\small\bfseries\sffamily, fill=white, inner sep=3pt},
    conn/.style={-{Stealth[length=2.5mm]}, thick},
    bconn/.style={{Stealth[length=2.5mm]}-{Stealth[length=2.5mm]}, thick},
    isoconn/.style={-{Stealth[length=2.5mm]}, very thick, densely dashed, red!70!black},
    pwrconn/.style={-{Stealth[length=2.5mm]}, very thick, black!75},
    lbl/.style={font=\scriptsize, midway, fill=white, inner sep=1.5pt},
  ]
  \node[fpga] (fpga) at (0,0) {FPGA\\{\small Actel A3P1000-PQG208M}\\[3pt]{\footnotesize\itshape MCS-48 Emulation\/}};

  \node[module, fill=gray!4] (pwr)   at (-7.5, 1.0) {Power Supply\\{\footnotesize \SI{3.3}{\volt} / \SI{1.5}{\volt}}};
  \node[module, fill=gray!4] (flash) at (-3.5, 3.5) {Flash ROM\\{\footnotesize S29AL016J}};
  \node[module, fill=gray!4] (xtal)  at (0, 3.5)    {Crystal\\{\footnotesize \SI{6}{\mega\hertz}}};
  \node[module, fill=gray!4] (nvram) at (3.5, 3.5)  {NVRAM\\{\footnotesize CY14B104NA}};
  \node[module, fill=gray!4] (bod)   at (-7.5, -1.3) {BOD\\{\footnotesize TPS3307}};
  \node[module, fill=gray!4] (jtag)  at (-7.5, 3.5)  {JTAG\\{\footnotesize Debug}};
  \node[module, fill=gray!4, minimum height=1.0cm] (lvl) at (5.6, -0.2) {Level Shift\\{\footnotesize CD4050}};
  \node[module, fill=gray!4, minimum height=1.0cm] (drv) at (5.6, -1.8) {Output Drivers\\{\footnotesize 54LS245-type}};
  \node[module, fill=gray!4] (freq)  at (-3.5, -3.5){Freq.\ Sweep\\{\footnotesize \SI{850}{\hertz}--\SI{1550}{\hertz}}};
  \node[module, fill=gray!4] (spi)   at (0, -3.5)   {SPI Master\\{\footnotesize (32-bit frames, 2023)}};
  \node[module, fill=gray!4] (ioexp) at (3.5, -3.5) {I/O \& Port\\{\footnotesize Expander}};

  \node[module, fill=green!8] (bcd)  at (9.3, 0.8)  {BCD Decoder\\{\footnotesize CD4511BE}};
  \node[module, fill=green!8] (disp) at (12.8, 0.8) {7-Segment\\{\footnotesize Displays ($\times 2$)}};

  \node[module, fill=blue!8] (wd)       at (-7.5, -6.5) {Watchdog\\{\footnotesize CD40106, CD4518}};
  \node[module, fill=blue!8] (btns)     at (-7.5, -8.5) {Buttons\\{\footnotesize Test, STW, L\"osch, T\"ur/V}};
  \node[module, fill=blue!8] (debounce) at (-3.5, -8.5) {Debounce};

  \node[module, fill=orange!10] (opto) at (3.5, -7)   {Optocouplers\\{\footnotesize SFH601-3 ($\times 4$)}};
  \node[module, fill=orange!10] (stm)  at (7, -7)     {STM32\\{\footnotesize F401RET6}};
  \node[module, fill=orange!10] (rtc)  at (10.8, -5.5)  {RTC\\{\footnotesize internal, \SI{32.768}{\kilo\hertz}}};
  \node[module, fill=orange!10] (sd)   at (10.8, -7)    {MicroSD\\{\footnotesize (FatFs)}};
  \node[module, fill=orange!10] (usb)  at (10.8, -8.5)  {USB\\{\footnotesize Type-C}};
  \node[module, fill=orange!10] (pwrin) at (0.6, -7) {D16, Y23 input\\{\footnotesize Base Board side}};

  \begin{scope}[on background layer]
    \node[pcb, draw=gray!50, fill=gray!2,
          fit=(fpga)(pwr)(flash)(xtal)(nvram)(bod)(jtag)(lvl)(drv)(freq)(spi)(ioexp)] (bb) {};
    \node[pcb, draw=green!40!black, fill=green!2, fit=(bcd)(disp)] (lb) {};
    \node[pcb, draw=blue!40!black, fill=blue!2, fit=(wd)(btns)(debounce)] (ub) {};
    \node[pcb, draw=orange!50!black, fill=orange!2, fit=(opto)(stm)(rtc)(sd)(usb)(pwrin)] (sb) {};
  \end{scope}

  \node[pcblabel, text=gray!60] at (bb.north) {PCB 1: Base Board};
  \node[pcblabel, text=green!40!black] at (lb.north) {PCB 4: LED Board};
  \node[pcblabel, text=blue!40!black] at (ub.south) {PCB 2: Upper Board};
  \node[pcblabel, text=orange!50!black] at (sb.south) {PCB 3: STM32 Diagnostics Board};

  \draw[conn] (pwr.east) -- node[lbl, above] {\SI{1.5}{\volt}} (pwr.east -| fpga.west);
  \draw[conn] (pwr.south) -- node[lbl, right] {\SI{3.3}{\volt}} (bod.north);
  \draw[conn] (bod.south) -- node[lbl, right] {Reset} (wd.north);
  \draw[bconn] (wd.east) -| (-5.45, -0.5) -- node[lbl, above] {Reset / !CS} (fpga.west |- 0, -0.5);
  \draw[conn] (jtag.east) -- ++(0.8, 0) |- (fpga.west |- 0, 1.6);
  \draw[conn] (lvl.west) -- node[lbl, above] {X4 inputs} (lvl.west -| fpga.east);
  \draw[conn] (drv.west -| fpga.east) -- (drv.west);
  \draw[conn] (drv.east) -- node[lbl, above] {X4 outputs} ++(1.9, 0);
  \draw[bconn] (flash.south) -- node[lbl, below left] {Firmware} ([xshift=-1.5cm]fpga.north);
  \draw[conn] (xtal.south) -- node[lbl, left] {CLK} (fpga.north);
  \draw[bconn] (nvram.south) -- node[lbl, below right] {Address / Data} ([xshift=1.8cm]fpga.north);
  \draw[conn] (fpga.240) -- (freq.north);
  \draw[conn] (fpga.south) -- (spi.north);
  \draw[bconn] ([xshift=1.8cm]fpga.south) -- (ioexp.north);
  \draw[conn] (fpga.east |- bcd.west) -- node[lbl, above, pos=0.3] {BCD} (bcd.west);
  \draw[conn] (bcd.east) -- (disp.west);
  \draw[conn] (btns.east) -- (debounce.west);
  \draw[conn] (debounce.north) |- (-5.05, -5.0) |- (fpga.west |- 0, -1.5);
  \draw[isoconn] (spi.south) -- ++(0,-1) -| node[lbl, right, text=red!70!black, font=\scriptsize\bfseries, pos=0.25] {Galvanic Isolation} (opto.north);
  \draw[conn] (opto.east) -- node[lbl, above] {SPI} (stm.west);
  \draw[bconn] (stm.north) |- (rtc.west);
  \draw[bconn] (stm.east) -- (sd.west);
  \draw[bconn] (stm.south) |- (usb.west);
  \coordinate (feedx) at (-0.3, 0);
  \coordinate (ledx) at (2.6, 0);
  \coordinate (branch) at (-0.3, -5.75);
  \draw[pwrconn] (bb.south -| feedx) -- node[lbl, left, pos=0.2] {\SI{5}{\volt}, GND} (pwrin.north -| feedx);
  \draw[pwrconn] (branch) -| node[lbl, below, pos=0.3] {GND} (opto.north -| ledx);
  \fill[black!75] (branch) circle (2pt);
\end{tikzpicture}%
}

	\caption[High-Level Block Diagram of the Retrofitted System]{High-level block diagram of the retrofitted test-diagnosis module. Four boards are shown: the Base Board with the \gls{fpga} core and the drivers of the outputs to the system connector \texttt{X4} (grey), the Upper Board with the operator controls and the watchdog (blue), the LED Board (green), and the galvanically isolated \gls{stm32} Diagnostics Board (orange). The dashed red arrow, labelled Galvanic Isolation, marks the \gls{spi} link from the \gls{spi} master to the optocouplers of the \gls{stm32} Diagnostics Board. The dark arrow carries the \SI{5}{\volt} supply and the ground of the Base Board through \texttt{X13}/\texttt{X12} to the parts of the \gls{stm32} Diagnostics Board on the Base Board side of the barrier; its branch, labelled GND, is the ground return of the optocoupler \glspl{led} (\Cref{sec:stm32_details}). Adapted from~\cite[Fig.~1]{torun2026}; the figure stays under the licence of the Working Paper, \mbox{CC~BY-NC-ND~4.0}.}
	\label{fig:overall_system}
\end{figure}

The retrofitted system comprises four boards (\Cref{fig:overall_system}) designed to collaboratively replicate and enhance the functionality of the original test-diagnosis module:

\begin{description}
	\item[Base Board (Unterplatine):] The central processing hub housing the Actel \texttt{A3P1000-\allowbreak PQG208M}~\gls{fpga} (\Cref{fig:schematic_sheet1}, designator \texttt{Y20} in~\Cref{tab:bom}) running the emulated~\gls{mcs48} core. The ProASIC3 family has a non-volatile, flash-based configuration memory that is insensitive to the configuration-memory upsets (firm errors) affecting \gls{sram}-based devices~\cite{guzman2011obsolescence}; radiation data for the family are reported in~\cite{poivey2011proasic3}. This board integrates the support circuitry, including power regulation (\texttt{Y1}, \texttt{Y2}), the \SI{6}{\mega\hertz} clock generation (\texttt{Q6}, \texttt{Y16}, \Cref{fig:schematic_sheet2}), \gls{nvram}~(\texttt{Y15}, \texttt{CY14B104NA}, \Cref{fig:schematic_sheet2}) for non-volatile storage of fault codes, the \gls{bod} supervisor (\texttt{Y3}, \Cref{fig:schematic_sheet1}), a~\gls{jtag} header (\texttt{X3}) for programming and debugging, and the necessary inter-board connectors (\texttt{X5}/\texttt{X6}, \texttt{X9}, \texttt{X13}). It also houses the interface to the system connector \texttt{X4}: CD4050 level shifters for its inputs (\texttt{Y7}, \texttt{Y9}--\texttt{Y11}, \texttt{Y17}) and 54LS245-type drivers for its outputs (\texttt{Y8}, \texttt{Y12}). The schematic is given in~\Cref{sec:app_schematic_sheets} and the layout in~\Cref{sec:app_base_board_schematics} (\Cref{fig:base_layout_top,fig:base_layout_bottom}). The Base Board layout shown corresponds to the inherited Rev~0.0 board, which was reused after rework (\Cref{subsec:hardware_challenges}).

	\item[Upper Board (Oberplatine):] Carries the operator interface and the supervision of the controller: the four legacy buttons (\texttt{S1}--\texttt{S4}) with their debouncing (\texttt{Y6}, \Cref{fig:schematic_sheet2}) and the external watchdog with its reset circuit (\texttt{Y4}, \texttt{Y5}, \texttt{Q1}, \Cref{fig:schematic_sheet1}). The buttons are those of the legacy front panel~\cite[Blatt~5/25]{rexroth1992wsp}: Test starts the self-test, \gls{stw} the fault query (\emph{Störstellen-Abfrage}), which shows the stored fault information entry by entry, and Lösch (clear) erases it. The button labelled \gls{tuerv} in the design files, V-SIGN on that panel, lets the speed-dependent signals be checked at standstill without the complete test sequence~\cite[Blatt~14/25, 19/25]{rexroth1992wsp}. It is connected to the Base Board by a ribbon cable between the headers \texttt{X8} and \texttt{X9}. The layout (Rev~1.0, designed by Leyer~\cite{leyer2023}) is shown in~\Cref{sec:app_upper_board_schematics} (\Cref{fig:upper_layout_top,fig:upper_layout_bottom}).

	\item[STM32 Diagnostics Board:] (\Cref{fig:schematic_sheet4}) This separate board provides enhanced diagnostic and logging capabilities, intentionally decoupled from the safety-critical anti-skid core logic using galvanic isolation. It features an \texttt{STM32F401RET6} microcontroller (\texttt{Y24}, Arm\textsuperscript{\textregistered}~Cortex\textsuperscript{\textregistered}-M4), chosen for its suitable peripheral set (\gls{spi}, \gls{sdio}, \gls{usb}, \gls{rtc}) and processing capability for the logging task, aligning with hardware used in the test fixture development~\cite{dolata2023}. Timestamps come from the microcontroller's internal~\gls{rtc}, clocked by a \SI{32.768}{\kilo\hertz} crystal (\texttt{Q8}) and backed by a CR2032 cell in holder \texttt{U1} (\texttt{BS-7}); an~\gls{sdio} interface to a MicroSD card (\texttt{J6}) provides persistent logging using FatFs, and a~\gls{usb} Type-C interface (\texttt{X15}) serves data retrieval. Communication with the Base Board~\gls{fpga} occurs via a galvanically isolated~\gls{spi} link (optocouplers \texttt{OK5}--\texttt{OK8}), which is meant to keep faults on the isolated side out of the \gls{fpga} domain (\Cref{sec:design_principles}); the parts of the board on the Base Board side of the barrier are not covered by it (\Cref{sec:stm32_details}). An isolated DC/DC converter (\texttt{Y23}; the footprint \texttt{Y22} is an unpopulated alternative) and a \SI{3.3}{\volt} low-dropout regulator (\texttt{Y19}, \texttt{TPS73633}) supply the isolated domain. The introduction of~\gls{usb} and~\gls{sd} interfaces necessitates future cybersecurity assessment (\Cref{sec:future_work}). See~\Cref{sec:app_stm32_board_schematics} (\Cref{fig:stm32_layout_top,fig:stm32_layout_bottom} for the Rev~1.0 layout).

	\item[LED Board:] Houses the two physical 7-segment~\gls{led} displays (\texttt{DS1}, \texttt{DS2}, Kingbright \texttt{SC03-12EWA}~\cite{kingbright2020}) showing the user the codes of \Cref{tab:requirement_basis} (e.g., 99 for no fault, 88 for the lamp test, 31 in the read-out of a sensor fault), driven by two driver ICs that decode~\gls{bcd} to 7-segment signals (\texttt{Y13}/\texttt{Y14}, \texttt{CD4511BE}). It plugs directly onto the Base Board (\texttt{X10}/\texttt{X11} into \texttt{X5}/\texttt{X6}). The schematic is part of \Cref{fig:schematic_sheet3}, and the layout (Rev~1.0) is shown in~\Cref{sec:app_led_board_schematics} (\Cref{fig:led_layout_top,fig:led_layout_bottom}).
\end{description}

In this report, \emph{the rebuilt unit} denotes the test-diagnosis module as rebuilt: the inherited Rev~0.0 Base Board with the Rev~1.0 Upper, LED and \gls{stm32} boards.

Physical interconnection utilises direct board-to-board mounting (LED Board onto Base Board via \texttt{X5}/\texttt{X6}; \gls{stm32} Diagnostics Board onto Base Board via \texttt{X13}) and cabled connections (Base Board to Upper Board via a ribbon cable between \texttt{X9} and \texttt{X8} for the button, interrupt and reset signals). This modular approach enhances testability and maintainability, fulfilling a key design principle (\Cref{sec:design_principles}). Leyer found the footprints of \texttt{X5} and \texttt{X6} on the Rev~0.0 Base Board wrong, because the connectors have pins of two sizes; on the prototype he bridged them with short wires, which he called an unstable solution, and he corrected the footprints in the Rev~1.0 layout~\cite[Sec.~9.2]{leyer2023}. The reused Base Board is Rev~0.0 and keeps these footprints; how the LED Board was connected to it on the rebuilt unit is not documented. Reusing previously modified hardware calls for an incoming inspection; apart from the inspection of the populated boards for the~\gls{bom}, none is documented for the rebuilt unit (\Cref{chap:safety_checklist}, item~CL-F2). Connector details are given in~\Cref{sec:app_connectors}.

\subsection{Outputs to the Vehicle}
\label{sec:vehicle_outputs}

The module acts on the vehicle only through the outputs of the system connector \texttt{X4}, which the 54LS245-type drivers \texttt{Y8} and \texttt{Y12} drive, some through series resistors and three through the transistors \texttt{Q3}--\texttt{Q5} (\Cref{fig:schematic_sheet2,fig:schematic_sheet3}). The fixture report names them: four frequency outputs to the axle cards (pins 12z, 10d, 12d and 12b), the outputs \enquote{AV-Sperre} A1 to A4 (16d, 14z, 14b, 14d) and \enquote{V-Test} (04d)~\cite[Sec.~5.1.5, 5.2]{dolata2023}. In the self-test, the module sends a frequency sweep to each of the four axle cards in turn~\cite[Sec.~5.1.5]{dolata2023}, which simulates a misbehaviour so that the axle cards' recognition of sliding can be tested~\cite[Sec.~2.9]{leyer2023}. According to the manual, the anti-skid control valves then act one after another, which is checked with a pressure gauge on the C pipe~\cite[Blatt~18/25]{rexroth1992wsp}. The check, according to the manual, runs automatically after the supply is switched on, or at standstill when the Test button is pressed with the braking signal present, and the spring-applied brake is to be applied for it~\cite[Blatt~14/25]{rexroth1992wsp}; the V-SIGN button is to be pressed only at standstill~\cite[Blatt~19/25]{rexroth1992wsp}. What the outputs AV-Sperre and V-Test act on in the vehicle is not described in these sources; the state of all outputs during power-up and during a watchdog or brown-out reset is not documented; and what prevents a self-test, with its sweep and valve actions, while the vehicle is moving is documented only as these instructions of the manual. Whether a reset by the watchdog or the \gls{bod} circuit in service leads the unchanged program into the automatic check that the manual describes after switching on, with its valve actions, has not been examined; it is an open question for a hazard analysis (\Cref{chap:safety_checklist}, item~CL-A4).

No safety function, no~\gls{sil} target and no safe state has been derived for the module (\Cref{chap:safety_checklist}, item~CL-A4). Where this report calls the emulated core \enquote{safety-related} or \enquote{safety-critical}, the words describe the intent of the design, not a classification.

% ------------------------------------------------------------
\section{FPGA Core Implementation Details}
\label{sec:fpga_core_details}

The core logic resides in the \texttt{A3P1000}~\gls{fpga} and is implemented in~\gls{vhdl}. The \gls{vhdl} sources (among them \texttt{atop.vhd}, \texttt{t48\_core.vhd}, \texttt{t8243\_sync\_notri.vhd}, \texttt{sweep\_top.vhd} and \texttt{Logger\_top.vhd}), the \gls{stm32} firmware and the fixture's test cases are not published with this report. \Cref{app:test_cases} gives only an adapted version of one \gls{yaml} test case of the fixture~\cite[Listing~3]{dolata2023}, in two steps, and reproduces one excerpt of twelve lines from \texttt{atop.vhd}.

\begin{description}
	\item[MCS-48 Emulation Core:] An emulation of the \gls{mcs48} processor serves as the central processing element. According to its original schematic, the legacy board carries a ROM-less member of the family (an \texttt{INS8035NI}; the schematic also names an NEC \texttt{D80C39HC}), which fetched its program from an external~\gls{eprom}. Following Leyer's approach~\cite{leyer2023}, an existing implementation (\texttt{t48\_core.vhd}) from the OpenCores repository\footnote{OpenCores website: \url{https://opencores.org/}. Validation of the specific core used relies on project testing.} was integrated unmodified. This core executes the original firmware, read from the Flash~\gls{rom} that replaced the~\gls{eprom}, preserving the legacy control logic and aiming at behavioural compatibility at the instruction level. The design history of the binary is not available. Its reuse rests on the service of the original module since the early 1990s (\Cref{chap:introduction}); no operational data are cited, and which evidence \mbox{\gls{iec61508}-3}~\cite{iec61508_3_2010} would require for reusing it has not been examined. The primary verification challenge lies in ensuring that this emulation is functionally and temporally equivalent to the original hardware within the target application's constraints (\Cref{subsec:equivalence_strategies,subsec:fpga_certification}).
	\item[Memory Interfaces:] The~\gls{fpga} implements the logic to interface directly with the external memory components located on the Base Board: an \texttt{S29AL016J} Flash~\gls{rom} (\texttt{Y21}) storing the executable firmware, and a \texttt{CY14B104NA}~\gls{nvram} (\texttt{Y15}) used for persistent storage of fault codes (\Cref{fig:schematic_sheet1,fig:schematic_sheet2}). The~\gls{fpga} generates the required address, data, and control signals (chip enable, output enable, write enable, byte enables) according to the timing requirements of these memory devices and the access patterns dictated by the emulated~\gls{mcs48} core. Checking the timing on these interfaces is part of Layer~1 verification (\Cref{sec:fpga_interface_testing}).
	\item[Peripheral Emulation and Interfaces:] The~\gls{fpga} logic replicates or interfaces with the essential peripherals of the legacy test-diagnosis board:
	      \begin{itemize}
		      \item \textit{7-Segment Display Driving Logic:} Generates the~\gls{bcd} signals and the latch enable (\texttt{FPGA\_7SEG\_LE}; at the drivers \texttt{CON\_7SEG\_LE}) required by the external \texttt{CD4511BE} display drivers located on the LED Board (\Cref{fig:schematic_sheet3}).
		      \item \textit{Input Button Handling:} The STW and Lösch buttons reach the emulated core directly on \texttt{T0} and \texttt{T1}; the Test and Tür/V buttons reach \texttt{P10}/\texttt{P11} and, combined on the Upper Board, the interrupt lines \texttt{!INT\_A}/\texttt{!INT\_B}, which the \gls{fpga} merges into the core's interrupt input~\cite{leyer2023} (\Cref{fig:schematic_sheet2}).
		      \item \textit{Frequency Sweep Generation:} A digital generator (\texttt{freq\_gen.vhd}, \path{sweep_top.vhd}) replaces the original analogue circuit based on a~\gls{vco}. It divides the \SI{6}{\mega\hertz} system clock by a factor taken from a lookup table,
		      \begin{equation}
			      f_\mathrm{out}(i) = \frac{f_\mathrm{clk}}{2 \cdot (N_\mathrm{div}(i) + 1)}, \qquad f_\mathrm{clk} = \SI{6}{\mega\hertz},
			      \label{eq:sweep}
		      \end{equation}
		      where an internal \SI{600}{\hertz} clock (\texttt{gen600Hz.vhd}) steps through the table and so produces a non-linear sweep profile. Each half period lasts $N_\mathrm{div} + 1$ cycles of the system clock, because the counter in \texttt{freq\_gen.vhd} counts from 0 up to $N_\mathrm{div}$ before the output toggles. $N_\mathrm{div}$ ranges from 1935 ($f_\mathrm{out} \approx \SI{1549.6}{\hertz}$) to 3529 ($f_\mathrm{out} \approx \SI{849.9}{\hertz}$); the table for the rising sweep has 33 entries and the table for the falling sweep 65 entries, so the sweeps cover \SI{850}{\hertz} to \SI{1550}{\hertz} and back, the values Leyer gives for the sweep of the original board~\cite[Sec.~2.9]{leyer2023}. The generator advances by one entry per period of the \SI{600}{\hertz} clock, which \texttt{gen600Hz.vhd} derives from the system clock by toggling its output every 5000 cycles, so the rising sweep reaches its end value after 31 steps, about \SI{52}{\milli\second}, and the falling sweep after 61 steps, about \SI{102}{\milli\second}; the tables end after 32 and 64 steps. Leyer reports that he checked the end frequencies, the duration and the profile of the sweep against the original board, without giving values~\cite[Sec.~4.10]{leyer2023} (\Cref{fig:sweep_comparison}); the captures read about \SI{920}{\hertz} at the low end of the original board, and the fixture report gives about \SI{180}{\milli\second} for the frequency change of the original test card (\Cref{subsec:fpga_interface_results}).
		      \item \textit{\gls{spi} Master Interface:} Implements an~\gls{spi} master (\texttt{Logger\_top.vhd}, \path{SPI_core.vhd}) that sends diagnostic frames to the~\gls{stm32} Diagnostics Board (\Cref{sec:stm32_details}). In Leyer's design of 2023, each frame has 32~bits: a 2-bit event-type prefix (error change or button change), the 6-bit button state and the 24-bit error vector. He implemented this logging function on the \gls{fpga} side but did not test it, and the \gls{stm32} had no code yet to interpret the frames~\cite[Sec.~3.3]{leyer2023}. The follow-on project of 2025 sends the 24-bit error vector as three bytes~\cite[pp.~20--21]{savchuk2025}. Which of the two loggers the \gls{fpga} design of the rebuilt unit contains is not documented (\Cref{sec:diagnostic_resolution}).
		      \item \textit{Watchdog/Reset Interaction:} The emulated controller services the external watchdog (\Cref{fig:schematic_sheet1}) through its chip-select output (\texttt{CS\_8048}, net \texttt{FPGA\_8048\_!CS} on \Cref{fig:schematic_sheet2}). Each rising edge of this line resets the watchdog counter: through the inverter \texttt{Y4C} and the capacitor \texttt{C8} it pulls the input of the Schmitt trigger \texttt{Y4B}, which \texttt{R3} and \texttt{R4} hold near the supply, briefly low, and \texttt{Y4B} drives the reset input of the counter \texttt{Y5} high (\Cref{fig:schematic_sheet1}). A falling edge produces no pulse. The \gls{fpga} receives the master reset (\texttt{FPGA\_8048\_!RESET}, also on \Cref{fig:schematic_sheet2}) from the watchdog and \gls{bod} circuit.
	      \end{itemize}
	\item[Clocking:] An external \SI{6}{\mega\hertz} crystal (\texttt{Q6} on \Cref{fig:schematic_sheet2}) provides the primary clock source for the~\gls{fpga} and the emulated~\gls{mcs48} core; the clock generation has been left as on the original board~\cite[Sec.~2]{leyer2023}.
\end{description}
The~\gls{vhdl} implementation aims for functional and timing equivalence, which the two-layer strategy described in~\Cref{chap:testing} is meant to verify; the state of the evidence is discussed in~\Cref{subsec:equivalence_discussion}.

% ------------------------------------------------------------
\section{STM32 Diagnostics Subsystem Details}
\label{sec:stm32_details}

The dedicated \gls{stm32} Diagnostics Board (\Cref{fig:schematic_sheet4}) provides enhanced monitoring capabilities, designed to be functionally independent from the safety-critical~\gls{fpga} core. Its purpose is logging and diagnostics; it is designed not to influence the primary anti-skid function except through the supply it shares with the Base Board (Isolated Power Supply, below).

\begin{description}
	\item[Microcontroller:] An \texttt{STM32F401RET6} (\texttt{Y24}, Arm\textsuperscript{\textregistered}~Cortex\textsuperscript{\textregistered}-M4) provides the processing power and peripheral set for the diagnostic functions. This microcontroller was chosen for its suitable performance, integrated peripherals (\gls{spi}, \gls{sdio}, \gls{usb}, \gls{rtc}), availability, and familiar development environment (STM32CubeIDE), and it matches the hardware used in the related test fixture project~\cite{dolata2023} (which used an \texttt{STM32F401RE} Nucleo board).
	\item[Galvanically Isolated SPI Slave:] Four optocouplers (\texttt{OK5}--\texttt{OK8}, \texttt{SFH601-3}, \Cref{fig:schematic_sheet4}) carry the three~\gls{spi} lines from the Base Board~\gls{fpga} (chip select, clock and data) and the reset line of the~\gls{stm32}; no signal runs in the opposite direction. Each stage inverts its signal, because a conducting \gls{led} turns on its phototransistor, which pulls the line held high by \texttt{R107}--\texttt{R110} low: the reset line, inverted before \texttt{OK7} by \texttt{Y4} (\Cref{fig:schematic_sheet1}), reaches the~\gls{stm32} with the right polarity, while the three~\gls{spi} lines arrive inverted (\Cref{sec:diagnostic_resolution}). This isolation is meant to keep faults on the isolated side of the non-safety-related \gls{stm32} Diagnostics Board from feeding back into the safety-related~\gls{fpga} domain through the signal path (\Cref{sec:design_principles}). The~\gls{stm32} acts as the~\gls{spi} slave, receiving the frames from the~\gls{fpga} master. Leyer's firmware (\texttt{USB\_RTC\_SD\_Test}~\cite{leyer2023}) receives them with the~\gls{spi} peripheral; the follow-on project replaced this by a software receiver that reads three bytes per event~\cite[p.~16]{savchuk2025}. The \gls{spi} core of 2023 (\texttt{SPI\_core.vhd}) toggles the \gls{spi} clock every 31 cycles of the \SI{6}{\mega\hertz} system clock (a comment in the source states 30), which gives
	      \begin{equation}
		      t_\mathrm{half} = \frac{31}{f_\mathrm{clk}} \approx \SI{5.2}{\micro\second}, \qquad f_\mathrm{SPI} = \frac{f_\mathrm{clk}}{62} \approx \SI{96.8}{\kilo\hertz}.
		      \label{eq:spi_clock}
	      \end{equation}
	      The typical turn-on and turn-off times of about \SI{3}{\micro\second} in the datasheet of the \texttt{SFH601}~\cite[p.~3]{vishay_sfh601} apply to a load of \SI{100}{\ohm} ($V_\mathrm{CC} = \SI{5}{\volt}$, $I_\mathrm{C} = \SI{2}{\milli\ampere}$). On this board the phototransistors have collector pull-ups of \SI{10}{\kilo\ohm} (\texttt{R107}--\texttt{R110}, \Cref{tab:bom}), and for that load the datasheet's typical curve of switching time against load resistance~\cite[Fig.~14]{vishay_sfh601} gives a turn-off time of roughly \SIrange{90}{100}{\micro\second}, read from the curve: many times the half-period of \Cref{eq:spi_clock}. On these figures the link has no timing margin at \SI{96.8}{\kilo\hertz}. The curve states no condition beyond an ambient temperature of \SI{25}{\celsius}, and the circuit differs from the datasheet's (a \SI{3.3}{\volt} pull-up supply, \SI{1}{\mega\ohm} base resistors \texttt{R99}--\texttt{R101} and \texttt{R106}), so this is an indication, not a measurement. The timing across the barrier has not been verified; the follow-on project reports only that an oscilloscope check of the \gls{stm32} Diagnostics Board showed no unstable communication~\cite[p.~22]{savchuk2025}. That the switching time contributed to the defects of the diagnostic path (\Cref{sec:diagnostic_resolution}) is a possibility, not a finding. On the Rev~0.0 prototype, Leyer tested the transfer of error codes over the isolated link with reused test programs and checked the codes in the log file~\cite[Sec.~4.11]{leyer2023}; the logging function of his final design was not part of this test (\Cref{sec:fpga_core_details}), and for the rebuilt unit see \Cref{tab:evidence_state}. The \gls{stm32} \gls{spi} peripheral supports only 8- and 16-bit data frames~\cite[Sec.~20.2.1]{st_rm0368}, so one 32-bit frame arrives as four bytes (or two 16-bit words) per transaction.
	\item[Data Logging:] An~\gls{sdio} interface connects to a standard MicroSD card (\texttt{J6}), enabling persistent storage of timestamped event logs. The~\gls{stm32} firmware uses the FatFs library (\texttt{ff.c}) to manage the file system on the~\gls{sd}~card, so that logs are written in a standard, accessible format (text files).
	\item[Real-Time Clock (RTC):] The microcontroller's internal~\gls{rtc}, which is clocked by the \SI{32.768}{\kilo\hertz} crystal \texttt{Q8} and backed by a CR2032 cell in holder \texttt{U1} (\texttt{BS-7}) on the \texttt{VBAT} pin, provides date and time stamps for logged events and is intended to keep time across power cycles. The time can be set via a \texttt{time.txt} file placed on the~\gls{sd}~card, as implemented by Leyer~\cite{leyer2023}.
	\item[USB Interface:] A~\gls{usb} Type-C connector (\texttt{X15}) gives a PC access to the~\gls{sd}~card as a Mass Storage Device, so that log files can be retrieved without removing the card; the firmware of the follow-on project also sends the error codes over a \gls{usb} virtual COM port to a host program~\cite[pp.~17, 19]{savchuk2025}. This interface introduces potential cybersecurity considerations (\Cref{sec:future_work}) that would need assessment before operational deployment.
	\item[Isolated Power Supply:] An isolated \SI{5}{\volt} DC/DC converter (\texttt{Y23}; \texttt{Y22} is an alternative footprint that is not populated) and the \SI{3.3}{\volt} regulator \texttt{Y19} create the power domain for the~\gls{stm32} and its peripherals, fed via connector \texttt{X12} from the Base Board, maintaining the galvanic isolation barrier. The solder jumper \texttt{SJ2} can bridge the two grounds across this barrier and must remain open in operation. The barrier runs through the \gls{stm32} Diagnostics Board, not around it. \texttt{X12} brings the \SI{5}{\volt} supply and the ground of the Base Board onto the board (the schematic notes that this ground, \texttt{GND3}, is connected to \texttt{GND}). On this side of the barrier lie the suppressor diode \texttt{D16} (\texttt{SMBJ5.0A}) across that supply, the input of \texttt{Y23}, and the \glspl{led} of the optocouplers with their series resistors \texttt{R82}--\texttt{R85}, which the logger and reset lines of the Base Board drive and which share only its ground (\Cref{fig:schematic_sheet4}). A short circuit of \texttt{D16} or of the converter input would load the \SI{5}{\volt} rail of the Base Board whether or not \texttt{SJ2} is open, so these two parts remain a fault path into the Base Board supply, and the \glspl{led} one into the lines that drive them. No test of the isolation is documented.
	\item[Status Indication:] An onboard~\gls{led} (\texttt{LED3}) provides basic status feedback for the diagnostics subsystem.
\end{description}
The C firmware was developed with STM32CubeIDE: Leyer's project \texttt{USB\_RTC\_SD\_Test} (\gls{spi} reception, \gls{rtc}, FatFs on the \gls{sd}~card, \gls{usb} mass storage)~\cite{leyer2023} and its revision by the follow-on project~\cite[pp.~16--19]{savchuk2025}.

% ------------------------------------------------------------
\section{PCB Design and Implementation}
\label{sec:pcb_design_impl}

The physical design of the Rev~1.0~\glspl{pcb}, prepared by Leyer~\cite{leyer2023}, focused on modularity and reliability, specifically addressing the~\gls{dfm} issues identified in the prior prototype (\Cref{subsec:hardware_challenges}).

\begin{description}
	\item[Modularity Benefits:] Separation onto four distinct boards (Base, Upper, STM32, LED) enhances electrical isolation between functional blocks (especially safety-related versus non-safety-related functions) and improves testability and maintainability during development and potentially in service (\Cref{sec:design_principles}).
	\item[Layout Practices:] Standard layout principles, as outlined by Jones~\cite{jones2004}, were applied, informed by the project's design guidelines. The boards use a four-layer stackup (top, power, ground, bottom) with ground planes (\Cref{sec:app_drc}) to ensure signal integrity and low-impedance return paths; \Cref{fig:base_layout_top} shows the top copper layer as an example. Decoupling capacitors were placed close to~\gls{ic} power pins, and trace routing followed common practice (e.g., avoiding sharp angles, using 45-degree bends, choosing widths by current and impedance, see~\cite[pp.~7, 10, 12]{jones2004}).
	\item[\gls{emc}/\gls{emi} Measures:] The design contains measures that bear on~\gls{emc}~\cite{ott2009emc}: the use of ground planes, widespread decoupling, physical separation of potentially noisy digital sections from sensitive analogue inputs, and the galvanic isolation of the~\gls{stm32} subsystem via optocouplers. Their effect has not been tested; the~\gls{en50121} series~\cite{en50121_3_2_2016} concerns, by its title, the~\gls{emc} of railway applications (\Cref{sec:validation_context}).
	\item[Addressing Rework and DFM Issues:] The Rev~1.0 layouts, prepared by Leyer at the end of his project~\cite[Sec.~9.2]{leyer2023}, correct the footprint and connectivity errors found on the Rev~0.0 prototype; this work replaced the displays and prepared the board files for manufacture. The inherited Base Board was not remade, so its \texttt{X5}/\texttt{X6} footprints are still those of Rev~0.0 (\Cref{sec:overall_architecture}). The assembly was guided by the Rev~1.0 schematic (\Cref{app:schematics}) and the Rev~1.0~\gls{bom}, which \Cref{app:bom_drc} continues; the notes there record open conflicts between the sources.
	\item[Connectors:] Standard, readily available headers and sockets (e.g., pin headers, ribbon cable connectors, TE~AMPMODU System~50 cited in~\Cref{sec:app_connectors}) were used for inter-board connections (e.g., \texttt{X9}, \texttt{X13}) to ensure component availability and ease of assembly and disassembly. The system connector to the vehicle side (\texttt{X4}) is a 48-pin DIN~41612 type~F connector (rows d, b, z).
\end{description}
The resulting hardware, documented in~\Cref{app:schematics}, consists of the Rev~1.0 Upper, LED and \gls{stm32} boards on the inherited Rev~0.0 Base Board (\Cref{sec:overall_architecture}).

% ------------------------------------------------------------
\section{Supervision and Isolation Features}
\label{sec:redundancy_failsafe}

The design includes features that supervise the emulated controller and its supply, regulate the supply and isolate the diagnostics subsystem (\Cref{sec:design_principles}). The watchdog and the~\gls{bod} circuit hold the emulated controller in reset when they detect a fault; that reset is supervision, not a derived safe state, and for the state of the outputs during it see \Cref{sec:vehicle_outputs}.

\begin{description}
	\item[External Hardware Watchdog Timer:] An independent hardware watchdog circuit on the Upper Board (\texttt{Y4}, a CD40106 oscillator, and \texttt{Y5}, a CD4518 counter, see~\Cref{fig:schematic_sheet1}) monitors the activity of the emulated~\gls{mcs48} core within the~\gls{fpga} via its chip-select line (\texttt{FPGA\_8048\_!CS}). If the line stops toggling for seven to eight periods of the watchdog clock (\Cref{subsec:fpga_interface_results}), which indicates a potential firmware hang or hardware fault, the watchdog pulls the reset line (\texttt{FPGA\_8048\_!RESET}) low through the transistor \texttt{Q1} and so holds the emulated controller in reset. Through the inverter \texttt{Y4} and the optocoupler \texttt{OK7}, the same reset also holds the microcontroller of the \gls{stm32} Diagnostics Board in reset (\Cref{fig:schematic_sheet1,fig:schematic_sheet4}). The reset does not reach the output drivers \texttt{Y8} and \texttt{Y12}, which are always enabled, so during a reset the outputs at \texttt{X4} follow what the~\gls{fpga} drives on its output pins (\Cref{sec:vehicle_outputs}). This provides a layer of supervision independent of the core processor emulation itself. Leyer tested its reaction on the Rev~0.0 prototype~\cite[Sec.~4.6]{leyer2023} (\Cref{fig:watchdog_reaction}); on the rebuilt unit, no observation of the watchdog after the re-soldering of the \gls{fpga} pins is documented (\Cref{subsec:fpga_interface_results}).
	\item[Brown-Out Detection (BOD):] A dedicated voltage supervisor~\gls{ic} (\texttt{TPS3307-18}, \texttt{Y3} on~\Cref{fig:schematic_sheet1}) monitors three supply rails: \SI{5}{\volt} through the divider \texttt{R103}/\texttt{R105} (\SI{2.4}{\kilo\ohm}/\SI{6.8}{\kilo\ohm}), \SI{3.3}{\volt} through \texttt{R102}/\texttt{R104} (\SI{4.7}{\kilo\ohm}/\SI{5.6}{\kilo\ohm}), and the \SI{1.5}{\volt} core rail directly. When a sensed voltage falls below the supervisor's threshold, the~\gls{bod} circuit pulls the same reset line low through \texttt{Q1} and holds the emulated controller in reset, which is meant to prevent unpredictable behaviour of the controller during power dips or brown-outs. With the typical thresholds of the \texttt{TPS3307-18}, \SI{2.93}{\volt}, \SI{1.68}{\volt} and \SI{1.25}{\volt}~\cite[pp.~2, 6]{ti_tps3307}, the design's trip points are about \SI{3.96}{\volt} on the \SI{5}{\volt} rail, \SI{3.09}{\volt} on the \SI{3.3}{\volt} rail and \SI{1.25}{\volt} (\SIrange{1.22}{1.29}{\volt}) on the core rail. They are trip points, not limits of safe operation. The trip point of \SI{3.09}{\volt} lies inside the range of \SIrange{3.0}{3.6}{\volt} given for the \SI{3.3}{\volt} I/O supply of the \gls{fpga}, but that of \SI{1.25}{\volt} lies below the \SIrange{1.425}{1.575}{\volt} specified for the core supply of the \texttt{A3P1000}~\cite[Table~7]{microsemi_ds0106}. Between about \SI{1.29}{\volt} and \SI{1.425}{\volt}, the core can therefore run outside its specified range without a reset. For the \SI{5}{\volt} rail, this work does not show that the parts it supplies work down to about \SI{3.96}{\volt}. This is a limitation of the design, recorded here and not measured. Leyer tested the circuit on the Rev~0.0 prototype and observed a reset when the supply fell below about \SI{4.1}{\volt}~\cite[Sec.~4.7]{leyer2023}; on the rebuilt unit, no test of it is documented (\Cref{sec:results_discussion_risk}).
	\item[Power Supply Regulation:] Dedicated low-dropout regulators (\texttt{Y1}, \texttt{TPS73615}, \SI{1.5}{\volt}; \texttt{Y2}, \texttt{TPS73633}, \SI{3.3}{\volt}) provide stable rails to the sensitive electronic components (\gls{fpga}, memories) from the potentially variable input supply voltage.
	\item[Galvanic Isolation:] As previously detailed (\Cref{sec:stm32_details}), the opto-isolated~\gls{spi} interface between the~\gls{fpga} (safety-related domain) and the~\gls{stm32} (non-safety-related domain) is designed as a barrier against faults feeding back from the isolated \gls{stm32} side into the \gls{fpga} domain. Its parts on the Base Board side remain a fault path into the Base Board (\Cref{sec:stm32_details}).
	\item[Built-In Self-Tests (Legacy):] While not a feature of the new design itself, the ability of the retrofitted system to execute the original system's self-tests (initiated via the \enquote{Test} button) is an important element for demonstrating functional equivalence. These legacy tests, executed by the emulated~\gls{mcs48} core within the~\gls{fpga}, provide a baseline check of core functionality, including the frequency sweep generator and display outputs. Checking them is part of the Layer~2 scenarios (\Cref{sec:system_scenario_testing}); for their execution on the rebuilt unit see \Cref{tab:evidence_state}. They contribute to fault detection, although they do not constitute a comprehensive~\gls{bist} with quantified diagnostic coverage.
\end{description}
These hardware features would form inputs to a safety argument for the retrofitted system; this report does not make one.
